Una espira es troba fixa en una regió on hi ha un camp magnètic uniforme en direcció perpendicular al full i cap endins, tal com s'indica a la figura de l'esquerra. En la figura de la dreta es mostra la gràfica de la variació del flux que travessa l'espira en funció del temps.

\figura[0.75]{fig1.png}

\begin{apartat}{1,25}
Determineu el sentit del corrent induït en l'interval de temps de 0~s a 4~s, en l'interval de 4~s a 8~s i en l'interval de 8~s a 12~s. Justifiqueu quina de les variacions de camp magnètic representades a sota (\textit{a} o \textit{b}) provoca la variació de flux.
\end{apartat}

\figura[0.75]{fig2.png}

\begin{apartat}{1,25}
Calculeu la intensitat de corrent elèctric en cada interval de temps si la resistència de l'espira és de $5\ \mathrm{m\Omega}$.
\end{apartat}

\nota{La llei d'Ohm estableix que $I = V/R$.}
